% GENERATED FILE. Edit canonical lessons, then run npm run generate:books.
% canonical-source-hash: fnv1a64:707675032fdb9cb5
% canonical-lessons: HI-C263-jhumka, HI-C263-bindi, HI-C263-curi, HI-C263-payal, HI-C263-basuri

\chapter{An earring, A bindi, A bangle, An anklet, A flute}
\label{ch:hi-an-earring-a-bindi-a-bangle-an-anklet-a-flute}
\hlchaptermodality{263}

\begin{chapteropening}
\textbf{By the end of this chapter:} \emph{I can say an earring, a bindi, a bangle, an anklet and a flute.}

Complete the last lesson of chapter 263: I can say an earring, a bindi, a bangle, an anklet and a flute.
\end{chapteropening}

\section[\texorpdfstring{\hi{झुमका} (jhumkā)}{झुमका (jhumkā)}]{\hi{झुमका} (jhumkā) — an earring}

\label{lesson:HI-C263-jhumka}



\begin{quote}
Before the new one: say the Hindi for a piece of jewellery, then the Hindi for a bangle.
\end{quote}

\subsection*{You'll want to know: \hi{झुमका}}
\textbf{\hi{झुमका}} — \emph{jhumkā} — \textquotedblleft{}an earring\textquotedblright{}.

\begin{cousinweb}
\textbf{\hi{झुमका}} (\emph{jhumkā}) is a tadbhava, a word inherited and reshaped by sound change. It comes from \emph{*jhumma}, \textquotedblleft{}crowd, cluster\textquotedblright{}, a form CDIAL reconstructs from its descendants (CDIAL 5404). CDIAL notes an added \emph{-k-}.
\end{cousinweb}

\begin{grammarlens}[title={What you've built}]
One more word for this chapter.
\end{grammarlens}

\subsection*{Your turn}
\begin{itemize}
\raggedright
  \item \emph{Say it:} \emph{jhumkā}
  \item \emph{Say it:} \emph{jhumkā}, once more
  \item \emph{Say it:} say \emph{jhumkā}
  \item \emph{Recall:} say the Hindi for a piece of jewellery, then the Hindi for a bangle, then say \emph{jhumkā} again
\end{itemize}

\begin{tcolorbox}[breakable,colback=teal!4,colframe=teal!35!black,title={Before you move on}]
What does \hi{झुमका} mean? (\textquotedblleft{}An earring\textquotedblright{}.) Say it once more.
\end{tcolorbox}
\section[\texorpdfstring{\hi{बिंदी} (bindī)}{बिंदी (bindī)}]{\hi{बिंदी} (bindī) — a bindi}

\label{lesson:HI-C263-bindi}



\begin{quote}
Before the new one: say the Hindi for a bangle, then the Hindi for an earring.
\end{quote}

\subsection*{You'll want to know: \hi{बिंदी}}
\textbf{\hi{बिंदी}} — \emph{bindī} — \textquotedblleft{}a bindi\textquotedblright{}.

\begin{cousinweb}
\textbf{\hi{बिंदी}} (\emph{bindī}) is built on Sanskrit \emph{bindu}, \textquotedblleft{}drop, spot\textquotedblright{}, with the ending \emph{-ī} (CDIAL 9240). CDIAL lists the inherited form of the word as \emph{būṁd}, \textquotedblleft{}drop\textquotedblright{}.
\end{cousinweb}

\begin{grammarlens}[title={What you've built}]
One more word for this chapter.
\end{grammarlens}

\subsection*{Your turn}
\begin{itemize}
\raggedright
  \item \emph{Say it:} \emph{bindī}
  \item \emph{Say it:} \emph{bindī}, once more
  \item \emph{Say it:} say \emph{bindī}
  \item \emph{Recall:} say the Hindi for a bangle, then the Hindi for an earring, then say \emph{bindī} again
\end{itemize}

\begin{tcolorbox}[breakable,colback=teal!4,colframe=teal!35!black,title={Before you move on}]
What does \hi{बिंदी} mean? (\textquotedblleft{}A bindi\textquotedblright{}.) Say it once more.
\end{tcolorbox}
\section[\texorpdfstring{\hi{चूड़ी} (cūṛī)}{चूड़ी (cūṛī)}]{\hi{चूड़ी} (cūṛī) — a bangle}

\label{lesson:HI-C263-curi}



\begin{quote}
Before the new one: say the Hindi for an earring, then the Hindi for a bindi.
\end{quote}

\subsection*{You'll want to know: \hi{चूड़ी}}
\textbf{\hi{चूड़ी}} — \emph{cūṛī} — \textquotedblleft{}a bangle\textquotedblright{}.

\begin{cousinweb}
\textbf{\hi{चूड़ी}} (\emph{cūṛī}) is a tadbhava, a word inherited from Sanskrit and reshaped by sound change. It comes from \emph{cūḍa}, \textquotedblleft{}bracelet\textquotedblright{} (CDIAL 4884). The Sanskrit word is known only from the old word-lists.
\end{cousinweb}

\begin{grammarlens}[title={What you've built}]
One more word for this chapter.
\end{grammarlens}

\subsection*{Your turn}
\begin{itemize}
\raggedright
  \item \emph{Say it:} \emph{cūṛī}
  \item \emph{Say it:} \emph{cūṛī}, once more
  \item \emph{Say it:} say \emph{cūṛī}
  \item \emph{Recall:} say the Hindi for an earring, then the Hindi for a bindi, then say \emph{cūṛī} again
\end{itemize}

\begin{tcolorbox}[breakable,colback=teal!4,colframe=teal!35!black,title={Before you move on}]
What does \hi{चूड़ी} mean? (\textquotedblleft{}A bangle\textquotedblright{}.) Say it once more.
\end{tcolorbox}
\section[\texorpdfstring{\hi{पायल} (pāyal)}{पायल (pāyal)}]{\hi{पायल} (pāyal) — an anklet}

\label{lesson:HI-C263-payal}



\begin{quote}
Before the new one: say the Hindi for a bindi, then the Hindi for a bangle.
\end{quote}

\subsection*{You'll want to know: \hi{पायल}}
\textbf{\hi{पायल}} — \emph{pāyal} — \textquotedblleft{}an anklet\textquotedblright{}.

\begin{cousinweb}
\textbf{\hi{पायल}} (\emph{pāyal}) goes back to a reconstructed \emph{*pādala}, \textquotedblleft{}belonging to the feet\textquotedblright{}, built on \emph{pāda}, \textquotedblleft{}foot\textquotedblright{} (CDIAL 8066), though CDIAL glosses the Hindi word only as \textquotedblleft{}sure-footed\textquotedblright{}.
\end{cousinweb}

\begin{grammarlens}[title={What you've built}]
One more word for this chapter.
\end{grammarlens}

\subsection*{Your turn}
\begin{itemize}
\raggedright
  \item \emph{Say it:} \emph{pāyal}
  \item \emph{Say it:} \emph{pāyal}, once more
  \item \emph{Say it:} say \emph{pāyal}
  \item \emph{Recall:} say the Hindi for a bindi, then the Hindi for a bangle, then say \emph{pāyal} again
\end{itemize}

\begin{tcolorbox}[breakable,colback=teal!4,colframe=teal!35!black,title={Before you move on}]
What does \hi{पायल} mean? (\textquotedblleft{}An anklet\textquotedblright{}.) Say it once more.
\end{tcolorbox}
\section[\texorpdfstring{\hi{बाँसुरी} (bā̃surī)}{बाँसुरी (bā̃surī)}]{\hi{बाँसुरी} (bā̃surī) — a flute}

\label{lesson:HI-C263-basuri}



\begin{quote}
Before the new one: say the Hindi for a bangle, then the Hindi for an anklet.
\end{quote}

\subsection*{You'll want to know: \hi{बाँसुरी}}
\textbf{\hi{बाँसुरी}} — \emph{bā̃surī} — \textquotedblleft{}a flute\textquotedblright{}.

\begin{cousinweb}
\textbf{\hi{बाँसुरी}} (\emph{bāṁsurī}) belongs with Hindi \emph{bãsī}, \textquotedblleft{}flute\textquotedblright{}, which CDIAL lists under Sanskrit \emph{vaṁśī}, \textquotedblleft{}flute, pipe\textquotedblright{}, built on \emph{vaṁśa}, \textquotedblleft{}bamboo\textquotedblright{} (CDIAL 11180). CDIAL does not list \emph{bāṁsurī} itself.
\end{cousinweb}

\begin{grammarlens}[title={What you've built}]
One more word for this chapter.
\end{grammarlens}

\subsection*{Your turn}
\begin{itemize}
\raggedright
  \item \emph{Say it:} \emph{bā̃surī}
  \item \emph{Say it:} \emph{bā̃surī}, once more
  \item \emph{Say it:} say \emph{bā̃surī}
  \item \emph{Recall:} say the Hindi for a bangle, then the Hindi for an anklet, then say \emph{bā̃surī} again
\end{itemize}

\begin{tcolorbox}[breakable,colback=teal!4,colframe=teal!35!black,title={Before you move on}]
What does \hi{बाँसुरी} mean? (\textquotedblleft{}A flute\textquotedblright{}.) Say it once more.
\end{tcolorbox}
